\documentclass[10pt, twoside]{article}

%handbook geometry
\usepackage[
    paperwidth=140mm,
    paperheight=100mm,
    margin=8mm,
    bindingoffset=4mm
]{geometry}

%packages
\usepackage[utf8]{inputenc}
\usepackage[italian]{babel} %To be modified in the other file
\usepackage{graphicx}
\usepackage{xcolor}
\usepackage{fancyhdr}
\usepackage{tcolorbox} 
\usepackage{wrapfig}

%font
\renewcommand{\familydefault}{\sfdefault}

%palette TBD

%pages and titles styles TBD

\begin{document}

\title{LANDER ADVANCE - MANUALE DI GIOCO}
\author{}
\date{}
\maketitle
\newpage

---------------------------------------------------------------------------------------------
\section{Protocollo XXX - Inizializzazione}
%prologo, il pilota lavora per una corporazione controllata da un pagliaccio

---------------------------------------------------------------------------------------------
\section{Intervista n° 339}
%discussione con il pagliaccio sugli obiettivi delle missioni; una specie aliena minaccia la Terra e solo facendo ridere i suoi membri, sparsi sui vari corpi celesti, la risparmieranno

---------------------------------------------------------------------------------------------
\section{Traiettoria di volo - Inizializzazione}
%descrizione del menù di selezione del target e delle difficoltà

---------------------------------------------------------------------------------------------
\subsection{Selezionare il corpo celeste}
%descrizione ed esempio sulla selezione del target

---------------------------------------------------------------------------------------------
\subsection{Scegliere l'area di atterraggio}
%descrizione dell'apertura della mappa di atterraggio e delle aree con relativi punteggi

---------------------------------------------------------------------------------------------
\section{Configurazione di volo - Inizializzazione}
%descrizione del menù di selezione lander e crew

---------------------------------------------------------------------------------------------
\subsection{Selezione del lander}
%descrizione ed esempio di selezione del lander, con statistiche 

---------------------------------------------------------------------------------------------
\subsection{Selezione delle risorse - FASE CRITICA}
\begin{tcolorbox}
ATTENZIONE: XXX non è resposabile della perdita di vite umane, il pilota lo è.
\end{tcolorbox}
%descrizione ed esempio di selezione della crew, con statistiche; ogni membro ha una massa diversa ma anche un moltiplicatore punti diverso, che indica quanto è bravo a far ridere (extra: gli ingegneri facilitano l'atterraggio ma fanno meno ridere, i comici lo rendono difficile ma danno più punti)

---------------------------------------------------------------------------------------------
\section{Inizializzazione della simulazione - Resti immobile}
%core gameplay: scopo, vittoria, sconfitta, punteggio finale

---------------------------------------------------------------------------------------------
\section{Impostazione della plancia}
%spiegazione dei comandi

---------------------------------------------------------------------------------------------
\section{Inizializzazione HUD}
%descrizione dell'HUD del pilota

---------------------------------------------------------------------------------------------------------------------
\section{Dichiarazione di limitazione delle responsabilità}
\begin{tcolorbox}
Accettando la missione e la destinazione, il pilota dichiara di essere consapevole che ogni corpo celeste comporta dei rischi differentii
\end{tcolorbox}
%descrizione di problematiche improvvise e risoluzione: perdita del sensore ottico

\end{document}